\clearpage
\appendix
\section*{Appendix}

\section{Results on UCFRep, Countix and QUVA}
\label{app:other}
The models of \cref{tab:main} are jointly trained on RepCount-A, UCFRep and
Countix; \cref{tab:app_other} reports their test results on the other two
benchmarks with the checkpoint selected on each benchmark's own validation split,
and their zero-shot results on QUVA (RepCount-A checkpoint and dataset token).
Three protocol differences matter when reading the UCFRep column. UCFRep has only
a training and an evaluation split; prior work reports the evaluation split and
does not describe a separate selection split, whereas we select on 90 training
videos held out by UCF101 group and train on the remaining 331 of 421. Context and
Sight\&Sound fine-tune their encoder on UCFRep, and CountLLM pretrains on
WebVid-10M and tunes a 7B language model; \method keeps its encoder frozen.

\begin{table}[h]
\centering
\caption{UCFRep (official evaluation split) and Countix (test); QUVA zero-shot.
$\dagger$: trained on RepCount-A, tested on UCFRep (cross-dataset).
$\ddagger$: released checkpoint evaluated by us. $\P$: Countix videos and MAE
normalisation not specified.}
\label{tab:app_other}
\small
\setlength{\tabcolsep}{3pt}
\begin{tabular}{lccccc}
\toprule
& \multicolumn{2}{c}{UCFRep} & \multicolumn{2}{c}{Countix} & QUVA\\
\cmidrule(lr){2-3}\cmidrule(lr){4-5}
Method & MAE & OBO & MAE & OBO & MAE\\
\midrule
RepNet~\cite{dwibedi2020repnet} & 0.474 & 0.371 & 0.364 & 0.697 & --\\
Context~\cite{zhang2020context} & 0.147 & 0.790 & -- & -- & --\\
Sight\&Sound~\cite{zhang2021sight} & 0.143 & 0.800 & 0.307 & 0.511 & --\\
TransRAC~\cite{hu2022transrac}$^\dagger$ & 0.640 & 0.324 & -- & -- & --\\
MFL~\cite{li2024mfl} & 0.388 & 0.510 & -- & -- & --\\
ESCounts~\cite{sinha2024escounts} & 0.216 & 0.704 & 0.276 & 0.673 & --\\
ESCounts$^\ddagger$ & -- & -- & -- & -- & 0.299\\
D$^2$-STX~\cite{wang2026d2stx} & 0.236 & 0.777 & 0.278 & 0.734 & --\\
CountLLM~\cite{yao2025countllm} & 0.116 & 0.839 & 0.215 & 0.801 & --\\
PAMS~\cite{gao2026pams} & 0.208 & 0.686 & 0.098$^\P$ & 0.879$^\P$ & --\\
\midrule
\method & \tabinput{tab/app_row}\\
\bottomrule
\end{tabular}
\end{table}

\Cref{tab:app_ablation} repeats the component ablations of \cref{tab:ablation} on
all three benchmarks.

% Component ablations on all three benchmarks (test). Rows written by temporac_cvpr/final_tables_frozen.py.
\begin{table*}[h]
\centering
\caption{Ablations on the three benchmarks (test, mean$\pm$std over seeds; each model scored once).}
\label{tab:app_ablation}
\small
\setlength{\tabcolsep}{4pt}
\begin{tabular}{lccccccc}
\toprule
& & \multicolumn{2}{c}{RepCount-A} & \multicolumn{2}{c}{UCFRep} & \multicolumn{2}{c}{Countix}\\
\cmidrule(lr){3-4}\cmidrule(lr){5-6}\cmidrule(lr){7-8}
Variant & seeds & MAE & OBO & MAE & OBO & MAE & OBO\\
\midrule
\tabinput{tab/app_ablation_rows}
\bottomrule
\end{tabular}
\end{table*}


\section{Count-aware encoder adaptation}
\label{app:adapt}
Counters that fine-tune their encoder are far stronger on UCFRep than those that
keep it frozen~\cite{zhang2020context,sinha2024escounts}, and in CountLLM adapting
the video encoder is the largest single gain~\cite{yao2025countllm}.
Back-propagating through the overlapping pyramid is too expensive, so we adapted
the encoder once, before extracting the pyramid: from every training video a
window of four non-overlapping 16-frame clips at a stride drawn from $\{2,4,8\}$,
rank-8 LoRA factors~\cite{hu2022lora} on the attention projections of all twelve
blocks and a small counting head that predicts the target density integral of
every temporal token; the epoch was selected on validation windows and the merged
encoder re-encoded every video. Adaptation lowered the validation MAE of the
two-view average from 0.190/0.258/0.144 (frozen encoder, three decoder seeds) to
0.175/0.255/0.134 (adapted, five seeds; RepCount-A, Countix, UCFRep), but on test
it raised the RepCount-A error and
helped only on Countix (\cref{tab:app_lm}, first two rows); we therefore keep the
encoder frozen.

\section{Language models}
\label{app:lm}
We studied four roles for a language model on top of \method; all use
Qwen3-VL-8B-Instruct~\cite{qwen3vl}. \Cref{tab:app_lm} collects their test
results; every model was scored once.

\paragraph{Counting directly.} Asked how often the action repeats, the frozen MLLM
names the action and describes one repetition correctly, but its counts are far
off: zero-shot MAE/OBO 0.548/0.31 on RepCount-A, 0.470/0.44 on UCFRep,
0.434/0.50 on Countix and 0.471/0.25 on QUVA.

\paragraph{Conditioning the decoder.} The MLLM's hidden states at the end of the
instruction and of its answer modulate the fused cells (FiLM~\cite{perez2018film}),
optionally with condition dropout. On validation this helped, on test it did not
(rows 3--4); a plausible reason is that the per-video state also encodes the
MLLM's own count guess and lets the decoder memorise training videos.

\paragraph{Refining the decoder.} Following CountLLM, a LoRA-tuned language model
(rank 16, 34M trainable parameters) reads the decoder's temporal tokens twice
after a description of the action, answers the count as text and feeds its hidden
states back into the density through a zero-initialised projection, with the
converged decoder frozen. It changed the validation error by $-0.9\%$ (two seeds)
and the test error little (row 5). Trained on the videos the decoder has already
fitted, the stage has almost no errors to learn from.

\paragraph{Choosing between tempo views.} The two views of
\cref{sec:shift} are right in different tempo ranges, and which one is right
depends strongly on the action: a gate that knows only the action class,
estimated for each validation video from the other validation videos of its
class, lowers the validation MAE of the two-view average by 4.6\%, 5.3\% and
6.8\% (RepCount-A, Countix, UCFRep). Without labels, we let the MLLM describe
the repeated action, embedded the description (BGE-base~\cite{xiao2023cpack}) and
weighted the views by the vote of the $n{=}32$ held-out videos with the most
similar descriptions, each voting for the view that was right for it, weighted by
how much the choice mattered. The memory holds held-out predictions, as in
blending~\cite{wolpert1992stacked}: the decoders' predictions for the validation
videos, with five-fold cross-validation over validation videos for every reported
validation number. Cross-validated, the memory lowered the validation MAE by
3.2\%, 4.0\% and 5.7\%; indexed by visual features instead of the description by
1.5\%, 2.1\% and 0.3\%; a single prior per benchmark changed it by $+0.9\%$,
$-0.6\%$ and $-4.6\%$; and a LoRA-tuned language model trained on the same
outcomes to answer which count is right by $-3.4\%$, $-3.8\%$ and $+2.1\%$. On
test none of them improved the plain average consistently (rows 6--10). Asking a
language model for the plausible cycle duration of the action instead is too
coarse (an error of a factor of 1.36 even when the class is known).

\begin{table}[h]
\centering
\caption{Test MAE of encoder adaptation and of the language-model roles (mean over
decoder seeds; number of seeds in brackets). Rows 6--10 reweight the two views of
the same five adapted decoders as row 2.}
\label{tab:app_lm}
\small
\setlength{\tabcolsep}{3pt}
\begin{tabular}{rlccc}
\toprule
& Variant & RC & UCF & CX\\
\midrule
1 & \method (frozen encoder) & \multicolumn{3}{c}{see \cref{tab:main,tab:app_other}}\\
2 & adapted encoder, two-view average (5) & 0.155 & 0.193 & 0.224\\
3 & frozen encoder + MLLM FiLM (5) & 0.142 & 0.206 & 0.238\\
4 & \quad + condition dropout (3) & 0.140 & 0.204 & 0.248\\
5 & adapted + LLM refining the decoder (2) & 0.153 & 0.207 & 0.230\\
6 & adapted + description memory (5) & 0.157 & 0.191 & 0.223\\
7 & adapted + visual-feature memory (5) & 0.159 & 0.207 & 0.223\\
8 & adapted + benchmark prior (5) & 0.156 & 0.193 & 0.223\\
9 & adapted + LoRA LLM selector (5) & 0.162 & 0.202 & 0.216\\
10 & adapted + out-of-fold memory (5) & 0.158 & 0.188 & 0.222\\
\bottomrule
\end{tabular}
\end{table}

\section{A slow-motion view}
\label{app:slow}
Shifting the pyramid one level \emph{down} is also exact: for $V''[2i]=V[i]$, a
stride-$s$ clip starting at an even frame reads $V''[j+sf]=V[j/2+(s/2)f]$, the
stride-$s/2$ clip of $V$, and only frames of $V$ are ever read, so an extra
stride-1 stream gives the $0.5\times$ view without interpolation. The decoder,
which never saw slowed-down training videos, over-counts it (validation MAE
0.263/0.415/0.211 against 0.184/0.287/0.128 on the original view of the adapted
decoders), and it is the best view only for the fastest Countix clips (cycles
under 0.4\,s).

\section{Training and validation errors differ on Countix}
\label{app:shift}
Decoders cross-fitted on five folds of the training union predict the original
view of held-out \emph{training} clips better than the $2\times$ view (MAE 0.232
vs.\ 0.257), whereas on validation clips the order is reversed (0.287 vs.\
0.254), mostly because slow validation clips are over-counted far more often. A
selector trained on out-of-fold training predictions therefore learns the wrong
preference for Countix (validation $+0.1\%$ on RepCount-A, $-0.2\%$ on Countix).
